% !TeX spellcheck = en_US
\documentclass[aspectratio=1610]{beamer}
\usetheme{Madrid}
\setbeamercovered{transparent}
\setbeamertemplate{navigation symbols}{}
\setbeamertemplate{itemize items}[circle]
\setbeamertemplate{enumerate items}[default]

\makeatletter
\define@key{beamerframe}{wide}[10pt]{%
	\def\beamer@cramped{\itemsep #1\topsep0.5pt\relax}}
\makeatother
	

\usepackage[english]{babel}
\usepackage[utf8]{inputenc}
\usepackage[T1]{fontenc}
\usepackage{array}
\usepackage{xcolor}
\usepackage{colortbl}
\usepackage{booktabs}

\title[Fall - Lecture 7] % (optional, use only with long paper titles)
{Intermediate Macroeconomics: Lecture  7 (Ch. 5)}
\subtitle{ECON 311 -- Intermediate Macroeconomics}

\author{Muhebullah Karimzada}
\institute{School of Economics \\\ University of Northern British Columbia}
\date{Fall 2026}


% ---------- UNBC COLORS ----------
\usepackage{xcolor}
\definecolor{UNBCGreen}{RGB}{0,102,51}
\definecolor{UNBCDark}{RGB}{0,74,37}
\definecolor{UNBCGold}{RGB}{189,155,96}
\definecolor{SoftCream}{RGB}{249,247,242}
\definecolor{SoftGray}{RGB}{244,246,248}
\definecolor{AccentBlue}{RGB}{41,98,255}
\definecolor{AccentOrange}{RGB}{230,126,34}
\definecolor{AccentTeal}{RGB}{0,150,136}
\definecolor{AccentRed}{RGB}{192,57,43}
\setbeamercolor{normal text}{fg=black,bg=white}
\setbeamercolor{structure}{fg=UNBCDark}
\setbeamercolor{title}{fg=white,bg=UNBCDark}
\setbeamercolor{subtitle}{fg=UNBCGold}
\setbeamercolor{frametitle}{fg=white,bg=UNBCDark}
\setbeamercolor{block title}{fg=white,bg=UNBCDark}
\setbeamercolor{block body}{fg=black,bg=SoftGray}
\setbeamercolor{block title alerted}{fg=white,bg=AccentRed}
\setbeamercolor{block body alerted}{fg=black,bg=SoftGray}
\setbeamercolor{item}{fg=UNBCDark}
\setbeamercolor{section in toc}{fg=UNBCDark}
\setbeamercolor{alerted text}{fg=AccentRed}
\setbeamercolor{author in head/foot}{fg=white,bg=UNBCDark}
\setbeamercolor{date in head/foot}{fg=white,bg=UNBCDark}
% ---------- UNBC BRANDING ----------
\usepackage{graphicx}
\usepackage{lmodern}
\usepackage{tcolorbox}
\usepackage{tikz}
\usetikzlibrary{positioning,calc}
\setbeamertemplate{headline}{}
\setbeamertemplate{footline}{%
  \leavevmode%
  \hbox{%
  \begin{beamercolorbox}[wd=.78\paperwidth,ht=2.8ex,dp=1.2ex,leftskip=1em]{author in head/foot}%
  \color{white}\textbf{ECON 311} \hspace{0.5em} \textcolor{UNBCGold}{|} \hspace{0.5em} Intermediate Macroeconomics%
  \end{beamercolorbox}%
  \begin{beamercolorbox}[wd=.22\paperwidth,ht=2.8ex,dp=1.2ex,rightskip=1em plus1fil]{date in head/foot}%
  \color{white}\insertframenumber/\inserttotalframenumber%
  \end{beamercolorbox}}%
}

\newtcolorbox{mybox}[2][]{
  colback=#2!8,
  colframe=#2,
  boxrule=0.8pt,
  arc=2mm,
  left=2mm,right=2mm,top=1.5mm,bottom=1.5mm,
  #1
}

\newcommand{\topbanner}[1]{%
\begin{tikzpicture}[remember picture,overlay]
\fill[UNBCGold] ([yshift=-0.65cm]current page.north west) rectangle ([yshift=-0.48cm]current page.north east);
\node[anchor=north west,font=\small\bfseries,text=UNBCDark] at ([xshift=0.35cm,yshift=-0.78cm]current page.north west) {#1};
\end{tikzpicture}%
}
\begin{document}

\begin{frame}[plain]
  \titlepage
\end{frame}

\begin{frame}[wide]{A Rise in the Depreciation Rate}
	\begin{columns} 
		% Column 1
		\begin{column}{.5\textwidth}
			\begin{itemize}
				\item Suppose the depreciation rate is higher rate ($\delta'>\delta$) due to an exogenous shock
				\vspace{1mm}
				\begin{itemize}
					\item The depreciation curve \textit{rotates} upward
					\item The investment curve remains unchanged
					\item The capital stock declines until it reaches the new steady state
					\vspace{1mm}
					\begin{itemize}
						\item This happens because depreciation exceeds investment ($sy_t>\delta y_t$)
					\end{itemize}
					\item The new steady state is located to the left, $k^{**}$
				\end{itemize}
			\end{itemize}
		\end{column}
		% Column 2    
		\begin{column}{.5\textwidth}			
			\begin{figure}
				\centering
				% [Figure not available] \includegraphics[width=1\linewidth]{"Figures/MACRO6_FIG05.06"}
			\end{figure}
		\end{column}%
	\end{columns}
\end{frame}

\begin{frame}[wide]{A Rise in the Depreciation Rate}
	\begin{itemize}
		\item What happens to output in response to this increase in the depreciation rate?
		\begin{itemize}
			\item The decline in capital reduces output
			\item Output declines rapidly at first, and then gradually settles down at its new, lower steady state level $y^{**}$
		\end{itemize}
	\end{itemize}
\end{frame}

\begin{frame}[wide]{Behaviour of Output after an Increase in $\delta$}
	\begin{figure}
		\centering
		% [Figure not available] \includegraphics[width=0.75\linewidth]{"Figures/MACRO6_FIG05.07"}
	\end{figure}
	
\end{frame}

\begin{frame}[wide]{The Principle of Transition Dynamics }
	\begin{itemize}
		\item If an economy is below steady state: 	It will grow
		\item If an economy is above steady state:	Its growth rate will be negative
		\item The farther below its steady state an economy is (in percentage terms), the faster the economy will grow
		\item The closer to its steady state, the slower the economy will grow
		\item This allows us to understand why economies grow at different rates
		\begin{itemize}
			\item Recall the convergence theory discussed earlier: poor countries tend to grow faster than developed countries
			\item Based on Solow model, it could mean the poor countries are just further away from its steady state
		\end{itemize}
		
	\end{itemize}
\end{frame}

\begin{frame}[wide]{Looking at Data through the Lens of the Solow Model}
	\begin{itemize}
			\item Rearrange the steady state $\bar{s}y^* = \delta k^*$:
			\[\frac{k^*}{y^*} = \frac{\bar{s}}{\delta}\]
			\item The capital to output ratio is the ratio of the investment rate to the depreciation rate
			\item Over the last 30 years, 
			\begin{itemize}
					\item The investment rate in South Korea has averaged more than 35 percent of GDP
					\item The investment rate in the United States has been about 25 percent
					\item In the poorest countries of the world, investment rates are only around 10 percent
				\end{itemize}
			\item However, it is conventional to assume that the depreciation rate is relatively similar across countries, about 7 percent
			\begin{itemize}
					\item if the $\delta$ is constant, then high investment rates should have high capital–output ratios
				\end{itemize}
		\end{itemize}
\end{frame}

\begin{frame}[wide]{Explaining Capital in the Solow Model}
	\begin{itemize}
			\item There is a strong positive relationship between the investment rate and capital-output ratio, as predicted in the model
		\end{itemize}
	\begin{figure}
			\centering
			% [Figure not available] \includegraphics[width=0.45\linewidth]{"Figures/MACRO6_FIG05.03"}
		\end{figure}
	
\end{frame}



\begin{frame}[wide]{Differences in Output per Capita}
	\begin{itemize}
			\item The Solow model shows TFP is more important in explaining per capita output
			\begin{itemize}
					\item Let $\alpha = 1/3$, $\frac{1}{1-1/3} = 3/2$ and $\frac{1/3}{1-1/3} = 1/2$
					\[y^* = \bar{A}^{\frac{1}{1-\alpha}} \left(\frac{\bar{s}}{\delta}\right)^{\frac{\alpha}{1-\alpha}}  = \bar{A}^{3/2}\left(\frac{\bar{s}}{\delta}\right)^{1/2}\]
					\item Differences in capital are explained in part by differences in investment rate, and in part by differences in the productivity parameter
				\end{itemize}
			\item We can calculate the significance of investment rates for explaining differences in $y^*$.	Income per capita ($y$) is decomposable into
			\begin{itemize}
					\item TFP differences
					\item Investment differences
				\end{itemize}
		\end{itemize}
\end{frame}

\begin{frame}[wide]{Differences in Output per Capita}
	\begin{itemize}
		\item Take the ratio of y* for two countries, assuming the depreciation rate is the same:
		\[\frac{y^*_{rich}}{y^*_{poor}} = \left(\frac{\bar{A}_{rich}}{\bar{A}_{poor}}\right)^{3/2} \times \left(\frac{\bar{s}_{rich}}{\bar{s}_{poor}}  \right)^{1/2}\]
		\item Five richest countries to GDP in the five poorest countries, they differ by a factor of 64 ($\frac{y^*_{rich}}{y^*_{poor}} = 64$ )
		\item In the richest countries, investment is about 25 to 30 percent. In the poorest countries, investment is about 7 percent. ($\left(\frac{\bar{s}_{rich}}{\bar{s}_{poor}}\right)^{1/2} = 4^{1/2} = 2$)
		\item Then, using the equation, $\left(\frac{\bar{A}_{rich}}{\bar{A}_{poor}}\right)^{3/2} = 32$
	\end{itemize}
\end{frame}



\begin{frame}[wide]{Understanding Differences in Growth Rates}
	\begin{itemize}
		\item Empirically, for Organization for Economic Cooperation and Development  (OECD) countries, transition dynamics holds:
		\vspace{1mm}
		\begin{itemize}
			\item Countries that were poor in 1960 grew quickly
			\item Countries that were relatively rich grew slower
		\end{itemize}
		\item For the world as a whole, on average, rich and poor countries grow at the same rate
		\vspace{1mm}
		\begin{itemize}
			\item Two implications of this:
			\begin{itemize}
				\item Most countries (rich and poor) have already reached their steady states
				\item Countries are poor not because of a bad shock, but because they have parameters that yield a lower steady state (determinants of the steady state invest rates and A)
			\end{itemize}
		\end{itemize}
	\end{itemize}
\end{frame}

\begin{frame}[wide]{Growth Rates in the OECD, 1960–2014}
	\begin{itemize}
		\item In this case, countries that are poor in 1960 would be far below their steady state
		\item Countries that are rich would be closer to (or even above) their steady state
	\end{itemize}
	\begin{figure}
		\centering
		% [Figure not available] \includegraphics[width=0.6\linewidth]{"Figures/MACRO6_FIG05.08"}
	\end{figure}
	
\end{frame}

\begin{frame}[wide]{Growth Rates around the World, 1960–2017}
	\begin{itemize}
		\item On average, rich and poor countries grow at the same rate
		\begin{itemize}
			\item But, the variance is high for countries with low initial GDP
		\end{itemize}
	\end{itemize}
	\begin{figure}
		\centering
		% [Figure not available] \includegraphics[width=0.6\linewidth]{"Figures/MACRO6_FIG05.09"}
	\end{figure}
	
\end{frame}

\begin{frame}[wide]{Case Study: South Korea and the Philippines}
	\begin{itemize}
		\item South Korea
		\begin{itemize}
			\item Grow at 6.4 percent per year
			\item Increased from 10 percent of U.S. income to 75 percent
		\end{itemize}
		\item Philippines
		\begin{itemize}
			\item Grow 2.5 percent per year
			\item Stayed around 10 percent of U.S. income
		\end{itemize}
		\item Transition dynamics predicts that
		\begin{itemize}
			\item South Korea was far below its steady state -- high growth rate
			\item Philippines is already at steady state -- low growth rate
		\end{itemize}
	\end{itemize}
\end{frame}

\begin{frame}[wide]{Case Study: South Korea and the Philippines}
	\begin{itemize}
		\item Assuming equal depreciation rates
		\[\frac{Y^*_{KR}}{Y^*_{US}} = \left(\frac{\bar{A}_{KR}}{\bar{A}_{US}}\right)^{3/2} \left(\frac{\bar{s}_{KR}}{\bar{s}_{US}}\right)^{1/2} \]
		
		\vspace{2mm}
		\item The long-run ratio of per capita incomes depends on
		\begin{itemize}
			\item the ratio of productivities (TFP levels)
			\item the ratio of investment rates
		\end{itemize}
	\end{itemize}
\end{frame}

\begin{frame}[wide]{Investment in South Korea and the Philippines, 1950–2017}
	\begin{itemize}
		\item The investment rate in South Korea rose sharply between 1960 and 1990, from about 15 percent in the 1950s to around 40 percent in the early 1990s
		\item The steady state income in Korea should increase according to the Solow model
		\begin{itemize}
			\item In order to transition to its new steady state, Korea grew at a faster rate	
		\end{itemize}
		\item Investment rates remained relatively stable in the United States and the Philippines
	\end{itemize}
	\begin{figure}
		\centering
		\includegraphics[width=0.55\linewidth]{"Figures/Investment in South Korea and the Philippines, 1950–2017"}
	\end{figure}
	
\end{frame}

\begin{frame}[wide]{Strengths and Weaknesses of the Solow Model}
	\begin{itemize}
		\item A key element of the Solow model is a production function that
		\begin{itemize}
			\item depends on capital and labor and an accumulation equation
			\item shows how forgoing consumption today leads to a higher capital stock tomorrow
		\end{itemize}
		\item The strengths of the Solow model:
		\begin{itemize}
			\item A baseline model that is able to determine how rich a country is in the long run (Steady State)
			\begin{itemize}
				\item Rich countries have a high rate of investment, a high TFP level, and/or a low rate of depreciation
			\end{itemize}
			\item The transition dynamics allows for an understanding of differences in growth rates across countries
			\item A country further from (closer to) the steady state will grow faster (slower)
			\begin{itemize}
				\item Matching the convergence pattern in the data for some countries
			\end{itemize}
		\end{itemize}
		
	\end{itemize}
\end{frame}

\begin{frame}[wide]{Strengths and Weaknesses of the Solow Model}
	\begin{itemize}
		\item The weaknesses of the Solow model:
		\begin{itemize}
			\item differences in investment rates explain only a small fraction of the differences in income across countries
			\item The much more important factor of TFP is still unexplained
		\end{itemize}
		\item It does not explain why different countries have different investment and productivity rates
		\begin{itemize}
			\item Economists extending the Solow model have endogenized the investment rate
			\item e.g., Korean government removed substantial barriers to investment, and the Philippines maintained these barriers
		\end{itemize}
		\item The model does not provide a theory of sustained long-run economic growth
		\begin{itemize}
			\item  Diminishing returns to capital accumulation in the production function mean that capital accumulation by itself cannot sustain growth
		\end{itemize}
	\end{itemize}
\end{frame}




%\begin{frame}[wide]{An Increase in the Investment Rate}
%	\begin{columns} 
%			% Column 1
%			\begin{column}{.5\textwidth}
%				\begin{itemize}
%					\item Suppose the investment rate increases permanently for exogenous reasons ($s'>s$)
%					\begin{itemize}
%						\item The investment curve rotates upward ($\bar{s}'\bar{A}  \bar{L}^{1-\alpha} K_t^\alpha > \bar{s}\bar{A}  \bar{L}^{1-\alpha} K_t^\alpha$)
%						\item The depreciation curve remains unchanged ($\delta K_t$)
%						\item The capital stock increases by transition dynamics to reach the new steady state
%						\begin{itemize}
%							\item This happens because investment exceeds depreciation ($sY_t>\delta K_t$)
%						\end{itemize}
%						\item The new steady state is located to the right
%					\end{itemize}
%				\end{itemize}
%			\end{column}
%			% Column 2    
%			\begin{column}{.5\textwidth}			
%				\begin{figure}
%					\centering
%					\includegraphics[width=1\linewidth]{"Figures/An Increase in the Investment Rate"}
%				\end{figure}
%			\end{column}%
%		\end{columns}
%\end{frame}

%\begin{frame}[wide]{Measuring the State of the Economy}
%	\begin{columns} 
%		% Column 1
%		\begin{column}{.5\textwidth}
%			\begin{table}[htbp]
%				\centering
%				\begin{tabular}{lrrrr}
%					& 2005  & 2006  & 2007  & 2008 \\
%					\hline
%					GDP   & 12.6  & 13.4  & 14.0    & 14.3 \\
%					(in trillions of \$) &       &       &       &  \\
%					Growth rate GDP &       & 0.063 & 0.045 & 0.021 \\
%					\hline
%				\end{tabular}%
%				\label{tab:addlabel}%
%			\end{table}%
%		\end{column}
%		% Column 2    
%		\begin{column}{.5\textwidth}			
%			
%			
%		\end{column}%
%	\end{columns}
%\end{frame}

\begin{frame}[wide]{Next Lecture}
	\begin{itemize}
		\item Next lecture will start Ch.6
	\end{itemize}
\end{frame}

\end{document}